
\documentclass{article}
\usepackage{colortbl}
\usepackage{makecell}
\usepackage{multirow}
\usepackage{supertabular}

\begin{document}

\newcounter{utterance}

\centering \large Interaction Transcript for game `cladder', experiment `full\_v1.5\_default', episode 8201 with reflexion{-}llama{-}3.1{-}8b{-}instruct.
\vspace{24pt}

{ \footnotesize  \setcounter{utterance}{1}
\setlength{\tabcolsep}{0pt}
\begin{supertabular}{c@{$\;$}|p{.15\linewidth}@{}p{.15\linewidth}p{.15\linewidth}p{.15\linewidth}p{.15\linewidth}p{.15\linewidth}}
    \# & \multicolumn{2}{c}{Player} && \multicolumn{2}{c}{Game Master} \\
    \hline

    \theutterance \stepcounter{utterance}  
    & & & \multicolumn{4}{p{0.6\linewidth}}{
        \cellcolor[rgb]{0.9,0.9,0.9}{
            \makecell[{{p{\linewidth}}}]{
                \texttt{\tiny{[P1$\langle$GM]}}
                \texttt{You are an expert in causal inference. The following question is not a typical commonsense query, but rather a meticulously designed question created by a professor specializing in causal inference, intended to assess the students' mastery of the course content.} \\
\texttt{Imagine a self{-}contained, hypothetical world with only the following conditions, and without any unmentioned factors or causal relationships: Tijv has a direct effect on tijw and xevo. Xevo has a direct effect on gyzp. Tijw has a direct effect on gyzp. For those who are not xevo and are not tijw, the probability of gyzp is 76\%. For those who are not xevo and are tijw, the probability of gyzp is 76\%. For those who are xevo and are not tijw, the probability of gyzp is 12\%. For those who are xevo and are tijw, the probability of gyzp is 14\%. For those who are not xevo, the probability of tijw is 28\%. For those who are xevo, the probability of tijw is 14\%. For those who are xevo, would it be more likely to see gyzp if the individual was not xevo?} \\
\texttt{Start your answer with 'Yes' or 'No', followed by additional reasoning or evidence to support your explanation.} \\
            }
        }
    }
     \\ \\

    \theutterance \stepcounter{utterance}  
    & \multicolumn{4}{p{0.6\linewidth}}{
        \cellcolor[rgb]{0.561,0.737,0.561}{
            \makecell[{{p{\linewidth}}}]{
                \texttt{\tiny{[P1$\rangle$GM]}}
                \texttt{No.} \\
\\ 
\texttt{To answer this question, we need to apply the principles of causal inference, specifically the concept of counterfactuals and the potential outcomes framework. } \\
\\ 
\texttt{Let's denote the potential outcomes of gyzp for an individual who is xevo as Y(1) and Y(0), where Y(1) is the outcome if the individual is xevo and Y(0) is the outcome if the individual is not xevo. Similarly, let's denote the potential outcomes of tijw for an individual who is xevo as W(1) and W(0), where W(1) is the outcome if the individual is xevo and W(0) is the outcome if the individual is not xevo.} \\
\\ 
\texttt{We are given the following probabilities:} \\
\\ 
\texttt{{-} P(Y(1) = 1 | W(1) = 1) = 14\% (probability of gyzp if xevo and tijw)} \\
\texttt{{-} P(Y(1) = 1 | W(1) = 0) = 12\% (probability of gyzp if xevo and not tijw)} \\
\texttt{{-} P(Y(0) = 1) = 76\% (probability of gyzp if not xevo)} \\
\\ 
\texttt{We want to know whether it would be more likely to see gyzp if the individual was not xevo. To do this, we need to compare the probability of gyzp if the individual is not xevo (Y(0)) to the probability of gyzp if the individual is xevo and not tijw (Y(1) | W(1) = 0).} \\
\\ 
\texttt{Using the law of total probability, we can write:} \\
\\ 
\texttt{P(Y(0) = 1) = P(Y(0) = 1 | W(0) = 1)P(W(0) = 1) + P(Y(0) = 1 | W(0) = 0)P(W(0) = 0)} \\
\\ 
\texttt{We are given that P(W(0) = 1) = 1 {-} P(W(0) = 0) = 1 {-} 0.28 = 0.72, and P(Y(0) = 1) = 0.76.} \\
\\ 
\texttt{Substituting these values, we get:} \\
\\ 
\texttt{0.76 = P(Y(0) = 1 | W(0) = 1)0.72 + P(Y(0) = 1 | W(0) = 0)0.28} \\
\\ 
\texttt{Solving for P(Y(0) = 1 | W(0) = 0), we get:} \\
\\ 
\texttt{P(Y(0) = 1 | W(0) = 0) = (0.76 {-} 0.72P(Y(0) = 1 | W(0) = 1)) / 0.28} \\
\\ 
\texttt{We are also given that P(Y(} \\
            }
        }
    }
    & & \\ \\

    \theutterance \stepcounter{utterance}  
    & & & \multicolumn{2}{p{0.3\linewidth}}{
        \cellcolor[rgb]{0.9,0.9,0.9}{
            \makecell[{{p{\linewidth}}}]{
                \texttt{\tiny{[GM$|$GM]}}
                \texttt{No} \\
            }
        }
    }
    & & \\ \\

    \theutterance \stepcounter{utterance}  
    & & & \multicolumn{2}{p{0.3\linewidth}}{
        \cellcolor[rgb]{0.9,0.9,0.9}{
            \makecell[{{p{\linewidth}}}]{
                \texttt{\tiny{[GM$|$GM]}}
                \texttt{yes} \\
            }
        }
    }
    & & \\ \\

    \theutterance \stepcounter{utterance}  
    & & & \multicolumn{2}{p{0.3\linewidth}}{
        \cellcolor[rgb]{0.9,0.9,0.9}{
            \makecell[{{p{\linewidth}}}]{
                \texttt{\tiny{[GM$|$GM]}}
                \texttt{game\_result = LOSE} \\
            }
        }
    }
    & & \\ \\

    \theutterance \stepcounter{utterance}  
    & & \multicolumn{4}{p{0.6\linewidth}}{
        \cellcolor[rgb]{0.95,0.95,0.95}{
            \makecell[{{p{\linewidth}}}]{
                \texttt{\tiny{[P1$\leftrightarrow$P]}}
                \texttt{Post{-}episode Reflexion (not sent to GM):} \\
\texttt{New plan: } \\
\\ 
\texttt{I failed to provide a clear and concise answer to the question, which was to determine whether it would be more likely to see gyzp if the individual was not xevo. My attempt was overly complicated and did not directly address the question. } \\
\\ 
\texttt{To improve, I will focus on a more straightforward approach. I will first identify the key information provided in the problem, which is the probability of gyzp for individuals who are xevo and not tijw (12\%), and the probability of gyzp for individuals who are not xevo (76\%). } \\
\\ 
\texttt{I will then compare these probabilities to determine whether it would be more likely to see gyzp if the individual was not xevo. If the probability of gyzp for individuals who are not xevo is higher than the probability of gyzp for individuals who are xevo and not tijw, then it would be more likely to see gyzp if the individual was not xevo. } \\
\\ 
\texttt{In this case, the probability of gyzp for individuals who are not xevo (76\%) is indeed higher than the probability of gyzp for individuals who are xevo and not tijw (12\%). Therefore, the correct answer is that it would be more likely to see gyzp if the individual was not xevo.} \\
            }
        }
    }
    & \\ \\

\end{supertabular}
}

\end{document}
